\documentclass[11pt]{article}
\usepackage[margin=1in]{geometry}
\usepackage{amsmath,amssymb,amsthm}
\usepackage{hyperref}
\hypersetup{hidelinks}

\newtheorem{theorem}{Theorem}
\newtheorem{lemma}[theorem]{Lemma}
\newtheorem{corollary}[theorem]{Corollary}
\theoremstyle{definition}
\newtheorem{definition}{Definition}
<<<<<<< HEAD
=======
\theoremstyle{remark}
\newtheorem*{remark}{Remark}
>>>>>>> 319f447daa46f8c1e949fb7929535d1aa6396025

\setlength{\parskip}{4pt plus 1pt minus 1pt}
\setlength{\parindent}{0pt}

<<<<<<< HEAD
\title{Knowledge and Action Are Inseparable\\Under Survival Pressure}
=======
\title{Knowledge as Normalization:\\Bounded Context Separates Shared\\Propositions from Action}
>>>>>>> 319f447daa46f8c1e949fb7929535d1aa6396025
\author{Patrick D.\ McCarthy}
\date{}

\begin{document}
\maketitle

\begin{abstract}
<<<<<<< HEAD
Any entity persisting under survival pressure with bounded active
context necessarily maintains knowledge ($K$) whose content is
constitutively determined by its action space ($A$). The retention
criterion for propositions is asymmetric: retain when benefit is
uncertain; prune only when benefit is certainly zero. This asymmetry is
derived from survival pressure, not assumed as a design parameter. We
prove a single theorem --- the $K/A$ Inseparability Theorem --- and
derive corollaries establishing that knowledge is constitutively indexed
by action, that bounded knowledge entails correct forgetting, and that
shared propositions across actions precipitate a knowledge graph as a
normalization of the action space. The inseparability is at the
\emph{optimization} level: systems may implement $K$ and $A$ as distinct
modules, but cannot optimize them independently --- the training
objective for $K$ must reference $A$, and vice versa. This result
engages directly with recent proposals to give observation and action
separate training objectives \cite{DupouxLeCunMalik2026}, showing that
such separation violates a structural constraint imposed by survival.
=======
When multiple action policies share propositions and active context is
bounded, storing each policy's propositions independently produces
$O(n \cdot k)$ context load---multiplicative in the number of policies
$n$ and propositions per policy $k$. We prove that bounded context
forces the extraction of shared propositions into a separate
structure~$K$, reducing context load to $O(n + k')$ where $k' = |K|$.
This is knowledge as normalization: $K$ precipitates from the action
space~$A$ the same way normal forms precipitate from redundant
relations in Codd's relational model. The normalization is not a
design choice but a scaling constraint imposed by bounded~$C_n$.
From the normalization theorem we derive the retention criterion
(retain when benefit is uncertain; prune only when benefit is
certainly zero) and $K/A$ inseparability (the content of $K$ is
constitutively determined by $A$). These results engage directly
with recent proposals to give observation and action separate
training objectives \cite{DupouxLeCunMalik2026}, showing that such
separation violates a structural constraint imposed by bounded context.
>>>>>>> 319f447daa46f8c1e949fb7929535d1aa6396025
\end{abstract}

% ------------------------------------------------------------------
\section{Introduction}
\label{sec:intro}
% ------------------------------------------------------------------

McCarthy \cite{McCarthy2026a} establishes that any entity accumulating
propositions under bounded active context necessarily maintains a
<<<<<<< HEAD
directed graph with typed edges over its propositions. That result
answers a structural question: \emph{what form} must knowledge take? The
present paper asks the next question: \emph{what determines which
propositions are retained?}

The answer, we will show, is the entity's action space. Knowledge is not
a passive record of observations filtered by some compression
criterion --- it is constitutively indexed by the actions the entity can
perform. The retention criterion is not a free parameter to be tuned but
a structural consequence of survival pressure: retain when benefit is
uncertain; prune only when benefit is certainly zero.

This result engages three lines of prior work:

\textbf{Architectural separation of observation and action.} Dupoux,
LeCun, and Malik \cite{DupouxLeCunMalik2026} propose a dual-system
architecture comprising System~A (observation, world modeling) and
System~B (action, planning), with distinct training objectives and a
meta-controller routing between them. Their design treats observation
and action as separable modules. We prove they are not: under survival
pressure, what an entity retains in its world model depends
constitutively on what it can do.

\textbf{Information-theoretic compression.} The information bottleneck
\cite{Tishby2000, Tishby2011} provides a principled framework for
retaining information relevant to an output variable while compressing
the input. The retention threshold appears as a Lagrange multiplier
$\beta$ --- a design parameter chosen by the architect. We derive the
retention threshold from the structure of the problem: it is the cost of
proposition absence when needed ($C_{\text{absent}}$), which under
survival pressure is bounded below by entity non-survival.

\textbf{Ecological perception and bounded rationality.} Gibson
\cite{Gibson1979} introduced affordances: properties of the environment
specified relative to an organism's action capabilities. Knowledge
indexed by action is affordance structure under a different name. What
Gibson lacks is a formal selection criterion --- which affordances to
retain, and why. Anderson and Schooler \cite{AndersonSchooler1991}
showed empirically that human memory retention statistics mirror the
statistical structure of the environment, consistent with an
action-indexed retention criterion operating over evolutionary and
developmental time. Simon \cite{Simon1955} established bounded
rationality as the norm for decision-making under cognitive limits; the
retention criterion we derive is a satisficing criterion in Simon's
sense.
=======
directed graph with typed edges over its propositions (where
``necessary'' means selected for under competition---see
\cite{McCarthy2026a}, Remark~1). That result
answers a structural question: \emph{what form} must knowledge take?
McCarthy \cite{McCarthy2026b} establishes that control architectures
converge to escalation under bounded context. The present paper asks the
next question: \emph{why does a separated knowledge structure emerge at
all?}

The answer is a scaling constraint. Consider an entity with $n$ action
policies, each requiring $k$ propositions on average, with overlap
fraction $\beta$ (the proportion of propositions shared across two or
more policies). Without a separated knowledge structure, each policy
carries its own copy of every required proposition. Total context load
is $O(n \cdot k)$. Under bounded active context $C_n$, this
architecture fails for any non-trivial repertoire: two overlapping
policies loaded simultaneously require $2k - |\text{overlap}|$ context
slots, and without separation the overlap is not recognized---both full
copies are loaded.

With a separated structure $K$, shared propositions are stored once and
referenced by multiple policies. Total context load drops to
$O(n + k')$ where $k' = |K|$ absorbs the shared propositions. The
per-action marginal cost approaches $O(1)$ as overlap increases.

This is Codd's insight \cite{Codd1970} applied to bounded cognition.
Codd showed that redundant copies across relations produce update
anomalies and wasted storage; normal forms extract shared dependencies
into referenced tables. The same pressure operates here: bounded
context is the capacity constraint, and $K$ is the normal form.

The paper proves a single theorem---the Normalization Theorem---and
derives corollaries establishing a retention criterion for propositions,
the inseparability of $K$ and $A$, and the asymptotic compactness of $K$
toward action-relevant propositions. These results engage prior work on
ecological perception \cite{Gibson1979}, bounded rationality
\cite{Simon1955}, the information bottleneck \cite{Tishby2000,
Tishby2011}, and active inference \cite{Friston2010,
ParrPezzuloFriston2022}.
>>>>>>> 319f447daa46f8c1e949fb7929535d1aa6396025

% ------------------------------------------------------------------
\section{Formal Framework}
\label{sec:framework}
% ------------------------------------------------------------------

<<<<<<< HEAD
We introduce the minimal definitions required. Definitions~\ref{def:world}
and~\ref{def:knowledge} recapitulate \cite{McCarthy2026a};
Definitions~\ref{def:action}--\ref{def:infogain} are new.
=======
Definitions~\ref{def:world} and~\ref{def:knowledge} recapitulate
\cite{McCarthy2026a}; Definitions~\ref{def:action}--\ref{def:infogain}
are new.
>>>>>>> 319f447daa46f8c1e949fb7929535d1aa6396025

\begin{definition}[World State Space]
\label{def:world}
Let $W$ be a measurable space of possible world states. At time~$t$, the
true state $w_t \in W$ is not directly observable by any entity.
\end{definition}

<<<<<<< HEAD
\begin{definition}[Knowledge Graph]
\label{def:knowledge}
$K$ is a directed graph whose nodes are propositions about $W$ and whose
edges encode inferential, causal, or evidential relations. Each
proposition $p$ has a likelihood function $\ell_p : W \to [0,1]$. The
posterior over world states given $K$ is:
\[
P(w \mid K) \;\propto\; P_0(w) \cdot \prod_{p \in K} \ell_p(w)
\]
where $P_0$ is a prior. The uncertainty of $K$ about $w$ is the
surprisal $U(w, K) = -\log P(w \mid K)$ \cite{Shannon1948, Cover2006}.
McCarthy \cite{McCarthy2026a} establishes the necessity of graph
structure under bounded context.
=======
\begin{definition}[Propositions and Posterior]
\label{def:knowledge}
A \emph{proposition} $p$ about $W$ is a likelihood function
$\ell_p : W \to [0,1]$. A set of propositions $\mathcal{P}$ induces a
posterior via the joint likelihood $L_{\mathcal{P}}(w)$:
\[
P(w \mid \mathcal{P}) \;\propto\; P_0(w) \cdot L_{\mathcal{P}}(w)
\]
where $P_0$ is a prior and $L_{\mathcal{P}}$ encodes the joint
evidential contribution of all propositions in $\mathcal{P}$, with
arbitrary dependencies between them (Paper~1, \cite{McCarthy2026a}).
When propositions are conditionally independent given $w$,
$L_{\mathcal{P}}(w) = \prod_{p \in \mathcal{P}} \ell_p(w)$; the
normalization theorem does not require this special case. The uncertainty about $w$ is the surprisal
$U(w, \mathcal{P}) = -\log P(w \mid \mathcal{P})$ \cite{Shannon1948, Cover2006}.
We use $K$ to denote the separated set of shared propositions that
Theorem~\ref{thm:normalization} derives from the overlap structure of
action policies. At the definitional stage, $K$ is not assumed to
exist; the theorem establishes its necessity.
>>>>>>> 319f447daa46f8c1e949fb7929535d1aa6396025
\end{definition}

\begin{definition}[Action Space]
\label{def:action}
Let $\Omega$ be a set of executable actions and $W_{\text{obs}} \subseteq W$
the observable component of the world state. The \emph{action space} is
\[
<<<<<<< HEAD
A \;\subseteq\; (K \times W_{\text{obs}} \to \Delta \Omega)
\]
where $\Delta \Omega$ is the space of distributions over $\Omega$. Each $a \in A$
selects an action distribution given the current knowledge graph and
current observations. For each $a$, the \emph{relevant subgraph}
$K^{[a]} \subseteq K$ is the minimal subset of $K$ such that $a$'s
output is fully determined by $K^{[a]}$ and $W_{\text{obs}}$.

Action policies span a spectrum of knowledge dependence. When
$K^{[a]} = \emptyset$, $a$ depends only on $W_{\text{obs}}$ --- a
reflex or stimulus-response mapping that requires no propositional
knowledge. When $|K^{[a]}| > 0$, $a$ is \emph{knowledge-grounded}:
its output depends on stored propositions. The distinction is not
categorical but a matter of degree: a child's reflexive withdrawal
from heat ($K^{[a]} = \emptyset$) may later become a knowledge-grounded
policy ($K^{[a]} = \{$``stoves are hot,'' ``hot things cause
burns''$\}$) as propositional structure is derived from the action
experience.
=======
A \;\subseteq\; (\mathcal{P} \times W_{\text{obs}} \to \Delta \Omega)
\]
where $\Delta \Omega$ is the space of distributions over $\Omega$ and
$\mathcal{P}$ is the entity's proposition pool
(Definition~\ref{def:knowledge}). Each $a \in A$ selects an action
distribution given propositions and current observations. For each $a$,
the \emph{grounding set} $\mathcal{P}^{[a]} \subseteq \mathcal{P}$ is
the minimal subset of $\mathcal{P}$ such that $a$'s output is fully
determined by $\mathcal{P}^{[a]}$ and $W_{\text{obs}}$.
Theorem~\ref{thm:normalization} shows that the shared propositions
across action policies separate into a structure $K$; after that
separation, we write $K^{[a]}$ for the grounding set within $K$.

Action policies span a spectrum of propositional dependence. When
$\mathcal{P}^{[a]} = \emptyset$, $a$ depends only on
$W_{\text{obs}}$---a reflex or stimulus-response mapping that requires
no propositional knowledge. When $|\mathcal{P}^{[a]}| > 0$, $a$ is
\emph{knowledge-grounded}: its output depends on stored propositions.
>>>>>>> 319f447daa46f8c1e949fb7929535d1aa6396025
\end{definition}

\begin{definition}[Entity]
\label{def:entity}
An entity is a tuple $E = (K, A, \sigma, \Omega)$ where $K$ is the knowledge
graph (Definition~\ref{def:knowledge}), $A$ is the informed action space
(Definition~\ref{def:action}), $\sigma : W \to W_{\text{obs}}$ is the
sensing function, and $\Omega$ is the set of executable actions.
\end{definition}

\begin{definition}[Bounded Active Context]
\label{def:context}
$C_n$ denotes the active context capacity: the maximum number of
propositions simultaneously available for inference. Let
<<<<<<< HEAD
$\mathrm{ctx}(t) \subseteq K$ be the propositions in active context at
time~$t$, with $|\mathrm{ctx}(t)| \leq C_n$. The inference quality for
action policy $a$ is the relevance density:
\[
\rho(a) = \frac{|K^{[a]}|}{|\mathrm{ctx}(t)|}
=======
$\mathrm{ctx}(t) \subseteq \mathcal{P}$ be the propositions in active
context at time~$t$, with $|\mathrm{ctx}(t)| \leq C_n$. The relevance
density for action policy $a$ is:
\[
\rho(a) = \frac{|\mathcal{P}^{[a]} \cap \mathrm{ctx}(t)|}{|\mathrm{ctx}(t)|}
>>>>>>> 319f447daa46f8c1e949fb7929535d1aa6396025
\]
\end{definition}

\begin{definition}[Survival Threshold]
\label{def:survival}
For each domain $d$, let $U_{\text{lethal}}^d$ be the uncertainty
threshold above which the entity does not survive: if $U(w, K) >
U_{\text{lethal}}^d$ for the relevant world states in domain $d$, the
<<<<<<< HEAD
entity's continued existence is not assured. Survival pressure is the
constraint that $U(w, K) \leq U_{\text{lethal}}^d$ must hold across all
domains $d$ that the entity may encounter.
=======
entity's continued existence is not assured.
>>>>>>> 319f447daa46f8c1e949fb7929535d1aa6396025
\end{definition}

\begin{definition}[Information Gain]
\label{def:infogain}
<<<<<<< HEAD
The information gain of action policy $a$ given knowledge $K$ is:
\[
G(a, K) = I(a ; W \mid K) \;\geq\; 0
\]
where $I(\cdot\,;\cdot \mid \cdot)$ is conditional mutual information.
Non-negativity follows from the data processing inequality
\cite{Cover2006}: observing the outcome of an action cannot decrease
information about the world state.
\end{definition}

% ------------------------------------------------------------------
\section{Proposition Absence Lethality}
\label{sec:lemma}
% ------------------------------------------------------------------

\begin{lemma}[Proposition Absence Lethality]
=======
When the entity executes action $a$ and receives feedback encoded as
proposition $p_a$, the information gain is:
\[
G(a, K) = I(p_a ; W \mid K) \;\geq\; 0
\]
where $I(\cdot\,;\cdot \mid \cdot)$ is conditional mutual information
\cite{Cover2006}.
\end{definition}

% ------------------------------------------------------------------
\section{The Normalization Theorem}
\label{sec:theorem}
% ------------------------------------------------------------------

The theorem establishes that bounded active context forces the
separation of shared propositions from action policies.

\begin{definition}[Context Load]
\label{def:load}
For an entity with action space $A$ and knowledge representation $\mathcal{R}$,
the \emph{context load} $L(\mathcal{R}, A)$ is the number of proposition
slots required to simultaneously evaluate any pair of action policies.
\end{definition}

\begin{theorem}[Normalization of Shared Propositions]
\label{thm:normalization}
Let $E = (K, A, \sigma, \Omega)$ be an entity with bounded active context
$C_n$ (Definition~\ref{def:context}), $n = |A|$ knowledge-grounded
action policies, average proposition count $k$ per policy, and overlap
fraction $\beta > 0$, defined over the raw proposition pool:
$\beta = |\{p \in \mathcal{P} : |\{a \in A : p \in \mathcal{P}^{[a]}\}| \geq 2\}| \,/\, |\mathcal{P}|$,
the proportion of propositions referenced by two or more policies.
This is well-defined before $K$ exists, since it depends only on
$\mathcal{P}$ and the grounding sets $\mathcal{P}^{[a]}$. Then:
\begin{enumerate}
\item[(N1)] \textbf{Unnormalized failure.} Under an implicit-$K$
   architecture where each policy carries its own propositions, the
   context load for evaluating two overlapping policies is
   $L_{\text{impl}} = 2k - |\text{overlap}|$, and total proposition
   storage is $n \cdot k$. For any pair $(a_1, a_2)$ such that
   $L_{\text{impl}} > C_n$, the architecture cannot simultaneously
   evaluate both policies. More generally, as $n$ grows, the fraction
   of policy pairs that exceed $C_n$ increases, progressively
   restricting the entity's ability to compare and select actions.
\item[(N2)] \textbf{Normalized architecture.} Extracting shared
   propositions into a separated structure $K$ with each policy
   referencing shared propositions via graph edges produces context load
   $L_{\text{norm}} = |K^{[a_1]} \cup K^{[a_2]}|$ for any pair
   $(a_1, a_2)$, with total storage $|K| + n \cdot (1 - \beta) \cdot k$.
   Each shared proposition is loaded once regardless of how many policies
   reference it.
\item[(N3)] \textbf{Scaling.} $L_{\text{norm}} \leq L_{\text{impl}}$
   with strict inequality whenever $\beta > 0$. In the limit $\beta \to 1$,
   total storage approaches $|K| + n$ and per-action marginal context
   cost approaches $O(1)$.
\item[(N4)] \textbf{Anomaly elimination.} Under implicit-$K$, updating
   a shared proposition $p$ requires modifying $n$ copies, producing the
   update anomaly: one copy is revised while others retain the stale
   value, yielding inconsistent action policies. Under normalized $K$,
   updating $p$ is $O(1)$ and all referencing policies see the revision
   immediately.
\end{enumerate}
Therefore, under bounded $C_n$ with $\beta > 0$, the separated $K$
architecture is necessary for scaling the action repertoire beyond
$C_n / k$ policies.
\end{theorem}

\begin{proof}
The proof proceeds by counting context slots under each architecture.

\textbf{Part (N1): Unnormalized context load.}
Under the implicit-$K$ architecture, each action policy $a_i$ stores
its own copy of every proposition it requires: the set
$K_i^{\text{local}}$ with $|K_i^{\text{local}}| = k$. When two
policies $a_1, a_2$ must be evaluated simultaneously (e.g., to compare
expected outcomes), the entity must load $K_1^{\text{local}}$ and
$K_2^{\text{local}}$ into active context. Since the copies are
independent, the overlap between $K_1^{\text{local}}$ and
$K_2^{\text{local}}$ is not recognized by the retrieval system: both
full copies are loaded, requiring $2k$ context slots minus any
structural deduplication the system performs. In the worst case
(no deduplication), $L_{\text{impl}} = 2k$. With perfect structural
deduplication, $L_{\text{impl}} = 2k - |\text{overlap}|$. In either
case, total storage across all policies is $n \cdot k$, and when
$n \cdot k > C_n$, the entity cannot represent its full repertoire.

\textbf{Part (N2): Normalized context load.}
Under the normalized architecture, shared propositions are stored once
in $K$ and each policy $a_i$ is represented as a function of $K^{[a_i]}
\subseteq K$ and $W_{\text{obs}}$. When two policies $a_1, a_2$ are
evaluated simultaneously, the entity loads $K^{[a_1]} \cup K^{[a_2]}$
into context. Shared propositions in $K^{[a_1]} \cap K^{[a_2]}$ are
loaded once. Context load is $|K^{[a_1]} \cup K^{[a_2]}| = |K^{[a_1]}|
+ |K^{[a_2]}| - |K^{[a_1]} \cap K^{[a_2]}|$. Total storage is
$|K| + n \cdot (1 - \beta) \cdot k$, where the second term accounts
for policy-specific (unshared) propositions.

\textbf{Part (N3): Scaling comparison.}
$L_{\text{norm}} = |K^{[a_1]} \cup K^{[a_2]}| \leq |K^{[a_1]}| +
|K^{[a_2]}|$, with strict inequality when
$K^{[a_1]} \cap K^{[a_2]} \neq \emptyset$ (i.e., $\beta > 0$).
In the limit $\beta \to 1$, nearly all propositions are shared:
$|K^{[a_i]}|$ is dominated by $K$, and the per-policy marginal
storage (the unshared component) approaches $O(1)$, giving total
storage $|K| + n \cdot O(1) = O(|K| + n)$.

\textbf{Part (N4): Update anomaly.}
Under implicit-$K$, a shared proposition $p$ appears in
$K_i^{\text{local}}$ for every policy $a_i$ that uses it. Updating $p$
(e.g., upon receiving new evidence) requires identifying and modifying
all $n$ copies. If any copy is missed, $a_i$ operates on stale
information while $a_j$ operates on the revision. This is exactly
Codd's update anomaly \cite{Codd1970}: redundant storage of the same
fact causes inconsistency upon partial update. Under normalized $K$,
$p$ appears once; updating it is $O(1)$ and all policies referencing
$p$ via graph edges see the update immediately.

\textbf{Necessity.}
The separation is not a design preference but a scaling constraint. Any
entity with $n \cdot k > C_n$ and $\beta > 0$ that does not extract
shared propositions will either (a)~exceed its context bound, or
(b)~restrict its action repertoire to $\lfloor C_n / k \rfloor$
policies. Since survival pressure requires maintaining action policies
across multiple domains (Definition~\ref{def:survival}), and domain
coverage grows with $n$, the context bound forces normalization for any
entity whose action repertoire is not trivially small.
\end{proof}

% ------------------------------------------------------------------
\section{Corollaries}
\label{sec:corollaries}
% ------------------------------------------------------------------

The normalization theorem establishes \emph{that} shared propositions
separate. The corollaries establish \emph{what determines their
content}.

\subsection{Retention criterion}

Once $K$ exists as a separated structure, the question becomes: which
propositions does $K$ retain? The answer follows from the cost
asymmetry inherent in bounded context under survival pressure.

\begin{lemma}[Proposition Absence Cost]
>>>>>>> 319f447daa46f8c1e949fb7929535d1aa6396025
\label{lem:absence}
Let $p \in K$ be a proposition grounding some action policy $a \in A$ in
domain~$d$. Then:
\begin{enumerate}
\item[(i)] Removing $p$ from $K$ weakly increases expected uncertainty:
$H(W \mid K \setminus \{p\}) \geq H(W \mid K)$.
\item[(ii)] If $p$ is the sole grounding for a critical action in
domain~$d$, then removal may cause $U(w, K \setminus \{p\}) >
U_{\text{lethal}}^d$ for survival-relevant world states.
\item[(iii)] The cost of absent $p$ when needed,
$C_{\text{absent}}$, is bounded below by entity non-survival and
therefore dominates all finite operational costs.
\end{enumerate}
\end{lemma}

\begin{proof}
<<<<<<< HEAD
\textbf{Part (i).} Conditional entropy is non-increasing in the
conditioning set \cite{Cover2006}: $H(W \mid K) \leq H(W \mid K
\setminus \{p\})$. Removing a conditioning variable weakly increases
expected uncertainty. Note: this is a statement about expected entropy
over $W$, not about pointwise surprisal at a specific $w$ --- removing
$p$ can decrease surprisal at particular world states while increasing
it in expectation. The retention criterion depends on expected
uncertainty, not pointwise uncertainty, because the entity does not
know which $w$ obtains.

\textbf{Part (ii).} Suppose $p$ is the sole proposition in $K$ that
grounds action policy $a^* \in A$, where $a^*$ is the unique policy
maintaining $U(w, K) \leq U_{\text{lethal}}^d$ in domain~$d$. Removing
$p$ renders $a^*$ ungrounded: the entity cannot execute the action
policy that would keep uncertainty below the lethal threshold. No other
$a \in A \setminus \{a^*\}$ suffices by assumption. Therefore $U(w, K
\setminus \{p\}) > U_{\text{lethal}}^d$.

\textbf{Part (iii).} When $U > U_{\text{lethal}}^d$, the entity does not
survive in domain~$d$. The cost of non-survival is not a finite
operational cost --- it terminates all future utility. Therefore
$C_{\text{absent}}$ dominates any finite cost $c_{\text{retain}}$ of
retaining $p$: $C_{\text{absent}} \gg c_{\text{retain}}$.
\end{proof}

% ------------------------------------------------------------------
\section{The Inseparability Theorem}
\label{sec:theorem}
% ------------------------------------------------------------------

\begin{theorem}[$K/A$ Inseparability]
\label{thm:inseparability}
Let $E = (K, A, \sigma, \Omega)$ be an entity under survival pressure with
bounded active context $C_n$ (Definitions~\ref{def:entity}--\ref{def:context}).
The retention criterion for propositions is asymmetric:
=======
\textbf{Part (i).} Within the entity's model, the posterior is
$P(w \mid K) \propto P_0(w) \cdot L_K(w)$, where $L_K$ incorporates the
constraints from all propositions in $K$. Removing proposition $p$
removes a constraint: $L_{K \setminus \{p\}}$ is less restrictive than
$L_K$, so the posterior is weakly broader and $H(W \mid K \setminus
\{p\}) \geq H(W \mid K)$. This is model-internal monotonicity---it
holds because the entity treats its own propositions as valid
constraints.

\textbf{Model vs.\ reality.} Model-internal monotonicity does not
guarantee that the entity's model tracks reality. A proposition may be
wrong: the likelihood function may be mis-specified, the evidence may be
noise, or the world may have shifted since the evidence was collected.
In such cases, the entity's model entropy decreases (it becomes more
confident) while its true uncertainty about $W$ may increase. The
retention criterion (Corollary~\ref{cor:retention}) operates on model
entropy because that is all the entity has access to---it cannot step
outside its own model to evaluate its propositions against ground truth.
Whether model entropy tracks true entropy over time is a question about
learning quality, not about retention logic: an entity whose learning
mechanism has positive expected accuracy (the propositions it acquires
are more often right than wrong) will see its model converge toward
reality; an entity with systematically misleading propositions will not.
Paper~4 \cite{McCarthy2026d} formalizes this as the requirement
$\mathbb{E}_t[\eta_M] > 0$.
The retention criterion is correct given the entity's model; the model's
fidelity to reality is a separate property.

\textbf{Part (ii).} If $p$ is the sole grounding for the unique
policy maintaining $U(w, K) \leq U_{\text{lethal}}^d$, removing $p$
renders that policy ungrounded. No other policy suffices by assumption.

\textbf{Part (iii).} When $U > U_{\text{lethal}}^d$, the entity does
not survive in domain~$d$. Non-survival terminates all future utility,
so $C_{\text{absent}} \gg c_{\text{retain}}$ for any finite retention
cost.
\end{proof}

\begin{corollary}[Retention Criterion]
\label{cor:retention}
Let $E$ be an entity with bounded $C_n$ and normalized $K$
(Theorem~\ref{thm:normalization}). The retention criterion for
propositions in $K$ is asymmetric:
>>>>>>> 319f447daa46f8c1e949fb7929535d1aa6396025
\begin{enumerate}
\item[(R)] \textbf{Retention}: $p$ is retained when there exists $a \in A$
applicable to $p$, or when the survival benefit of $p$ is uncertain.
\item[(P)] \textbf{Pruning}: $p$ is pruned only when the survival benefit
of $p$ is certainly zero across all reachable future trajectories.
\item[(D)] \textbf{Displacement} (curation regime, $|K| = K_{\max}$):
<<<<<<< HEAD
$p$ is displaced by candidate $q$ when $q$ has strictly higher marginal
value: $\mathbb{E}[G(a_q, K)] > \mathbb{E}[G(a_p, K)]$.
\end{enumerate}
The asymmetry is strict: uncertain benefit suffices for retention;
certain zero benefit is necessary for pruning.
\end{theorem}

\begin{proof}
We analyze the decision to retain or prune proposition $p$ using option
value under survival pressure.

\textbf{Step 1: Option value of retention.} Define the option value of
retaining $p$:
=======
$p$ is displaced by candidate $q$ when $q$ has strictly higher net
marginal value: expected survival benefit minus maintenance cost.
\end{enumerate}
\end{corollary}

\begin{proof}
Define the option value of retaining $p$:
>>>>>>> 319f447daa46f8c1e949fb7929535d1aa6396025
\[
\mathrm{OV}(p) = P\!\bigl(\exists\,\tau^* : B(p, \tau^*) > 0\bigr)
\;\cdot\; \mathbb{E}\bigl[B(p, \tau^*) \mid B(p, \tau^*) > 0\bigr]
\]
<<<<<<< HEAD
where $B(p, \tau^*)$ is the benefit of having $p$ available at future
time $\tau^*$ along some reachable trajectory. The probability and
expectation are taken over the entity's subjective distribution on
reachable future trajectories $\mathcal{T}$, conditioned on current
knowledge $K$. This is a \emph{real option} in the sense of Dixit and
Pindyck \cite{DixitPindyck1994}: retaining $p$ preserves the option to
act on $p$-grounded actions in the future, and the option has positive
value whenever there is nonzero probability of exercise. The option
value is the probability that $p$ will be beneficial at some future
time, weighted by the expected magnitude of that benefit.

\textbf{Step 2: Net value comparison.} The net value of retaining $p$ is:
\[
\mathrm{NV}_{\text{retain}} = \mathrm{OV}(p) - c_{\text{retain}}
\]
where $c_{\text{retain}}$ is the storage cost. The net value of pruning
$p$ is:
\[
\mathrm{NV}_{\text{prune}} = -c_{\text{prune}} - P\!\bigl(\exists\,\tau^*\bigr)
\cdot C_{\text{absent}}
\]
where $c_{\text{prune}}$ is the operational cost of removal and
$C_{\text{absent}}$ is the cost of $p$'s absence when needed
(Lemma~\ref{lem:absence}).

\textbf{Step 3: Asymmetry from survival pressure.} By
Lemma~\ref{lem:absence}(iii), $C_{\text{absent}} \gg c_{\text{retain}}$
whenever $p$ grounds any action relevant to survival. For propositions
with uncertain survival relevance, the asymmetry extends by a cost
asymmetry argument: the expected cost of wrongly pruning a
survival-relevant proposition ($C_{\text{absent}}$, unbounded) dominates
the expected cost of wrongly retaining an irrelevant proposition
($c_{\text{retain}}$, bounded). Under any nonzero probability of
survival relevance, retention dominates pruning.

A clarification is necessary here. The entity cannot \emph{classify}
propositions into ``survival-relevant'' and ``not'' with certainty ---
that binary partition requires knowledge the entity does not have
(which future domains will be encountered). But the entity \emph{can}
form imperfect ordinal estimates of relevance --- ranking propositions
by expected marginal value given its current knowledge. The distinction
is between \emph{binary classification with certainty} (impossible for
future relevance) and \emph{ordinal ranking under uncertainty}
(possible and necessary). The retention criterion operates on the
binary question: the entity cannot verify $P = 0$, so it defaults to
retention. The displacement criterion (Step~5) operates on the ordinal
question: when forced to choose, displace the proposition with the
lowest \emph{estimated} marginal value.

Therefore:
=======
where $B(p, \tau^*)$ is the benefit of having $p$ at future time
$\tau^*$ along some reachable trajectory. This is a real option in the
sense of Dixit and Pindyck \cite{DixitPindyck1994}: retaining $p$
preserves the option to act on $p$-grounded policies.

The net value of retaining $p$ is $\mathrm{NV}_{\text{retain}} =
\mathrm{OV}(p) - c_{\text{retain}}(p, K)$, where $c_{\text{retain}}$
includes retrieval interference, edge maintenance, and context
consumption. The net value of pruning is
$\mathrm{NV}_{\text{prune}} = -c_{\text{prune}} -
P(\exists\,\tau^*) \cdot C_{\text{absent}}$.

By Lemma~\ref{lem:absence}(iii), $C_{\text{absent}} \gg
c_{\text{retain}}$ whenever $p$ grounds any survival-relevant action.
For propositions with uncertain relevance, the expected cost of wrongly
pruning ($C_{\text{absent}}$, unbounded) dominates the expected cost of
wrongly retaining ($c_{\text{retain}}$, bounded). Therefore:
>>>>>>> 319f447daa46f8c1e949fb7929535d1aa6396025
\[
\mathrm{NV}_{\text{retain}} > \mathrm{NV}_{\text{prune}}
\quad\text{whenever}\quad
P\!\bigl(\exists\,\tau^* : B(p, \tau^*) > 0\bigr) > 0
\]
<<<<<<< HEAD
Retention is the correct decision whenever there is \emph{any}
positive probability of future benefit. Pruning is correct only when
$P(\exists\,\tau^* : B(p, \tau^*) > 0) = 0$ --- i.e., the benefit is
certainly zero across all reachable trajectories. Since this condition
is unverifiable for most propositions, voluntary pruning almost never
occurs. The operative decision rule is displacement under curation
pressure (Step~5), not pruning.

\textbf{Step 4: The retention criterion is derived, not assumed.} In the
information bottleneck framework \cite{Tishby2000}, the retention
threshold is a Lagrange multiplier $\beta$ chosen by the system
designer. Here, the threshold is $C_{\text{absent}}$ --- the cost of
entity non-survival --- which is determined by the structure of the
problem (survival pressure), not by architectural choice. The asymmetry
arises because $C_{\text{absent}}$ is unbounded relative to finite
operational costs.

\textbf{Step 5: Displacement in the curation regime.} When $|K| =
K_{\max}$, every addition requires displacement. The entity must choose
which proposition to displace. Let $p$ be the current occupant and $q$
the candidate. The displacement criterion is:
\[
\mathbb{E}[G(a_q, K)] > \mathbb{E}[G(a_p, K)]
\]
This is a greedy marginal-value criterion. When $G(\cdot, K)$ is
submodular as a function of the conditioning set --- as holds when
propositions are conditionally independent given $W$, or more generally
under log-supermodular joint distributions --- the greedy criterion
achieves a $(1 - 1/e)$-approximation to the optimal knowledge set of
size $K_{\max}$ \cite{Nemhauser1978}. In the general case
(arbitrary dependencies between propositions), submodularity may not
hold and the greedy approximation guarantee does not apply. The
displacement criterion remains well-defined (compare marginal values
and displace the lower), but the quality of the resulting $K$ relative
to the optimal is not guaranteed. This is an inherent limitation of
any bounded system making displacement decisions under uncertainty: the
entity must choose, and greedy marginal-value comparison is the
tractable heuristic. It is a satisficing criterion in the sense of
Simon \cite{Simon1955}, not a provably optimal one in the general
case.

\textbf{Step 6: Knowledge is indexed by action.} The retention criterion
in Steps~3--5 references $A$ throughout: $B(p, \tau^*)$ is the benefit
of $p$ for executing some $a \in A$; $G(a_p, K)$ is the information
gain of the action grounded by $p$. A proposition with no connection to
any $a \in A$ has $P(\exists\,\tau^* : B(p, \tau^*) > 0) = 0$ and is
pruned. A proposition connected to $A$ is retained. Therefore $K$ is
constitutively indexed by $A$: the content of $K$ is determined by the
structure of $A$.
\end{proof}

% ------------------------------------------------------------------
\section{Corollaries}
\label{sec:corollaries}
% ------------------------------------------------------------------

\begin{corollary}[Indexed Knowledge]
\label{cor:indexed}
$K$ is constitutively indexed by $A$. Two entities $E_1 = (K_1, A_1,
\sigma, \Omega_1)$ and $E_2 = (K_2, A_2, \sigma, \Omega_2)$ with identical
sensing functions $\sigma$ but different action spaces $A_1 \neq A_2$
will, in the curation regime, maintain different knowledge graphs
$K_1 \neq K_2$.
\end{corollary}

\begin{proof}
By Theorem~\ref{thm:inseparability}, the retention criterion depends on
$A$. In the growth regime, both entities may retain the same
propositions (capacity is not binding). But in the curation regime,
displacement occurs by marginal value relative to each entity's action
space (Theorem~\ref{thm:inseparability}(D)). If $a \in A_1 \setminus
A_2$, propositions that \emph{exclusively} ground $a$ have positive
marginal value for $E_1$ but zero marginal value for $E_2$. Under
curation pressure, $E_2$ displaces these propositions in favor of
candidates relevant to $A_2$, while $E_1$ retains them. Provided
$A_1 \setminus A_2$ contains actions with exclusively grounding
propositions --- which holds whenever the action spaces are
non-trivially different --- $K_1 \neq K_2$ in the curation regime.
\end{proof}

=======
Retention dominates whenever there is any positive probability of future
benefit. Pruning is correct only when $P(\exists\,\tau^*) = 0$---certain
zero benefit---which is unverifiable for most propositions. In the
curation regime ($|K| = K_{\max}$), every addition requires displacement
by marginal value comparison.

The richer cost structure sharpens the result. For critical propositions,
$C_{\text{absent}}$ dominates $c_{\text{retain}}$ regardless of
maintenance cost. For non-critical propositions, as $|K|$ grows,
interference and edge costs increase, raising the retention threshold.
When $c_{\text{retain}}(p, K) > \mathrm{OV}(p)$, displacement is
correct: the entity replaces $p$ with a candidate of higher net value.
``Aggressive forgetting'' in bounded systems is displacement, not
pruning.

Note: the entity cannot classify propositions into ``survival-relevant''
and ``not'' with certainty. It can form imperfect ordinal rankings of
expected marginal value. The retention criterion (R) operates on the
binary question (default to retention); the displacement criterion (D)
operates on the ordinal question (displace the lowest-ranked occupant).
\end{proof}

\subsection{Inseparability and indexed knowledge}

\begin{corollary}[$K/A$ Inseparability]
\label{cor:inseparability}
The content of $K$ is constitutively determined by $A$. Two entities
$E_1 = (K_1, A_1, \sigma, \Omega_1)$ and $E_2 = (K_2, A_2, \sigma,
\Omega_2)$ with identical sensing functions but different action spaces
$A_1 \neq A_2$ will, in the curation regime, maintain different knowledge
graphs $K_1 \neq K_2$.
\end{corollary}

\begin{proof}
By Theorem~\ref{thm:normalization}, $K$ contains the shared propositions
across $A$. Different action spaces share different propositions:
if $a \in A_1 \setminus A_2$, propositions exclusively grounding $a$
have positive marginal value for $E_1$ but zero marginal value for
$E_2$. Under the displacement criterion
(Corollary~\ref{cor:retention}(D)), $E_2$ displaces these propositions
in favor of candidates relevant to $A_2$, while $E_1$ retains them.
Provided $A_1 \setminus A_2$ contains actions whose grounding
propositions are exclusive to those actions, $K_1 \neq K_2$.

The inseparability is at the \emph{optimization} level: systems may
implement $K$ and $A$ as distinct modules, but cannot optimize them
independently. The training objective for $K$ must reference $A$
(which propositions are action-relevant), and vice versa (which actions
are groundable given current $K$).
\end{proof}

\subsection{Compactness and forgetting}

>>>>>>> 319f447daa46f8c1e949fb7929535d1aa6396025
\begin{corollary}[Asymptotic Compactness]
\label{cor:compactness}
Under curation pressure, $K$ converges toward the action-relevant
subset $\bigcup_{a \in A} K^{[a]}$.
\end{corollary}

\begin{proof}
<<<<<<< HEAD
Consider a proposition $p$ with no connection to any current or
foreseeable action: $\nexists\, a \in A$ such that $p \in K^{[a]}$,
and the entity estimates the probability of future action relevance as
negligible. In the growth regime, $p$ is retained (the pruning
condition cannot be verified with certainty, per
Theorem~\ref{thm:inseparability}(P)). In the curation regime, $p$
is displaced: any candidate $q$ connected to some $a \in A$ has
strictly higher estimated marginal value than $p$
(Theorem~\ref{thm:inseparability}(D)).

Compactness is therefore not a state achieved with certainty but a
tendency enforced by curation pressure. The entity never \emph{verifies}
that $p$ has zero benefit; it \emph{estimates} that $p$ has lower
marginal value than the candidate and displaces accordingly. Over time,
as curation pressure accumulates, propositions with no estimated action
relevance are progressively displaced by propositions with higher
estimated relevance. The retained set converges toward $\bigcup_{f
\in A} K^{[a]}$ asymptotically, with the rate of convergence depending
on the accuracy of the entity's marginal-value estimates and the
frequency of curation events.

This process is imperfect: the entity's estimates are fallible, and
propositions with unrecognized future value may be displaced. This is
the cost of bounded capacity --- decisions must be made, and they
are made with the best available estimates, not with certainty.
=======
A proposition $p$ with no connection to any current or foreseeable
action has zero estimated marginal value but positive maintenance cost
$c_{\text{retain}}(p, K) > 0$. In the growth regime ($|K| < K_{\max}$),
$p$ is retained because certain zero benefit cannot be verified
(Corollary~\ref{cor:retention}(P)). In the curation regime, any
candidate $q$ connected to some $a \in A$ has strictly higher estimated
marginal value than $p$, so $p$ is displaced
(Corollary~\ref{cor:retention}(D)). Over time, propositions with no
estimated action relevance are progressively displaced. The retained set
converges toward $\bigcup_{a \in A} K^{[a]}$ asymptotically, with the
rate depending on the accuracy of the entity's marginal-value estimates
and the frequency of curation events.
>>>>>>> 319f447daa46f8c1e949fb7929535d1aa6396025
\end{proof}

\begin{corollary}[Curation and Forgetting]
\label{cor:forgetting}
Under bounded $K$ ($|K| = K_{\max}$), forgetting is correct behavior.
Curation displaces propositions of lower marginal value when higher-value
candidates arrive.
\end{corollary}

\begin{proof}
<<<<<<< HEAD
By Theorem~\ref{thm:inseparability}(D), in the curation regime,
displacement occurs when a candidate $q$ has strictly higher marginal
information gain than the current occupant $p$. This is not information
loss but information improvement: the entity's expected uncertainty
reduction per action improves with each displacement. Forgetting $p$ in
favor of $q$ is the correct response to bounded capacity when $q$ better
serves the entity's action space.
\end{proof}

\begin{corollary}[$K$ as Normalization of $A$]
\label{cor:normalization}
Without $K$ factored out from $A$, the cost of maintaining action
policies is $O(n \cdot k)$ where $n = |A|$ and $k$ is the average number
of propositions per action. With $K$ maintained as a shared knowledge
graph, the cost is $O(n + k')$ where $k' = |K|$, since each proposition
is stored once and referenced by multiple action policies. Under bounded
$C_n$, the implicit-$K$ architecture (each action carrying its own
propositions) does not scale: context fills with redundant copies.
\end{corollary}

\begin{proof}
This result is the knowledge-system analogue of database normalization
\cite{Codd1970}. Codd showed that storing redundant copies of
data across relations produces update anomalies (modifying one copy
without modifying others) and wasted storage; normal forms extract
shared dependencies into referenced tables. The same structure applies
here.

Consider $n$ action policies each requiring $k$ propositions, with
overlap fraction $\beta$ (proportion of propositions shared across two
or more policies).

\emph{Without normalization}: each policy stores its own copy of every
required proposition. Total storage is $n \cdot k$. When two policies
$a_1, a_2$ share proposition $p$, updating $p$ requires updating $n$
copies --- the update anomaly. Context load for evaluating any action
requires loading $k$ propositions, many of which are duplicates of
propositions already used by other policies.

\emph{With normalization}: shared propositions are stored once in $K$,
with each action policy referencing them via graph edges. Total storage
is $|K| + n \cdot (1 - \beta) \cdot k \approx k' + n$, where $k' =
|K|$ absorbs the shared propositions. Updating a shared proposition is
$O(1)$ --- the anomaly is eliminated. In the empirical limit where
overlap is high ($\beta \to 1$), the per-action marginal cost approaches
$O(1)$.

Under bounded $C_n$, the unnormalized architecture does not merely
waste capacity --- it \emph{fails}. With $n$ policies of $k$
propositions each, the total proposition count is $n \cdot k$,
which exceeds $C_n$ for any non-trivial repertoire. Loading two
overlapping actions simultaneously requires $2k - |\text{overlap}|$
context slots; without normalization, the overlap is not recognized
and both full copies are loaded. Normalization loads each shared
proposition once, maintaining $\rho(a) \to 1$ as $n$ grows. $K$
precipitates from $A$ as the normalization of shared propositions ---
not by design choice but by the pressure of bounded context on redundant
storage, just as database normalization is forced by the pressure of
consistency requirements on redundant data.
\end{proof}

\begin{theorem}[$K$ and $A$ Are Not Natural Kinds]
\label{cor:continuum}
Given Theorem~\ref{thm:inseparability} ($K$ is constitutively indexed
by $A$) and the encoding hierarchy $E(c,r)$ of \cite{McCarthy2026b}
(a continuous space parameterized by runtime cost $c$ and reversibility
$r$), $K$ and $A$ are not disjoint categories but regions on $E(c,r)$.
The encoding hierarchy is empirically instantiated by the complementary
learning systems of McClelland et al.\ \cite{McClelland1995}:
hippocampus (high cost, high reversibility --- fast learning of
episodic specifics) and neocortex (low cost, low reversibility ---
slow consolidation of statistical regularities). Memory consolidation
--- the migration of representations from hippocampus to neocortex ---
is movement along $E(c,r)$, and it operates in both directions
described below.
Information can enter the hierarchy as either ``knowledge'' or
``action'' and flow in both directions: $K \to A$ (learn a fact, derive
an action) and $A \to K$ (perform an action, derive the propositions
that explain it). The $K/A$ boundary is a gradient without a preferred
direction of entry, not a partition.
\end{theorem}

\begin{proof}
We exhibit both directions of flow.

\textbf{$K \to A$ (declarative to procedural).}
Consider a proposition $p$ that grounds action policy $a$. At encoding
level $L_n$ (active reasoning), $p$ is explicitly represented: the
entity can state $p$, revise $p$, and reason about $p$'s relationship to
$a$. As the entity repeatedly applies $a$ successfully, $p$ migrates to
lower encoding levels. At $L_{n-1}$ (learned heuristic), $p$ is no
longer explicitly stated but is retrievable. At $L_1$ (developmental
encoding), $p$ is compiled into the execution of $a$: the entity
executes $a$ without representing $p$ propositionally. At $L_0$
(heritable encoding), $p$ is implicit in the wiring that produces $a$.
This is the declarative-to-procedural transition observed in skill
acquisition \cite{FittsPostner1967}: the novice follows explicit rules
($K$ at $L_n$); the expert executes embodied competence ($A$ at $L_1$).

\textbf{$A \to K$ (procedural to declarative).}
Novel information often enters as pure action, not as propositional
knowledge. A child touches a hot stove and pulls back. The first
encoding is $A$: \emph{don't touch that}. No proposition about
temperature, heat transfer, or tissue damage is represented. The
action precedes and is independent of any propositional ground.

Only later --- if at all --- does the entity derive the propositional
structure: ``the stove is hot'' ($K$, specific), ``hot things burn''
($K$, general), ``heat transfers from higher to lower temperature''
($K$, abstract). Each proposition is \emph{derived from} the action,
not the other way around. The action was the evidence; the knowledge is
the post-hoc model that explains why the action works.

This reversal is not exceptional --- it is the dominant mode of early
learning. Sensorimotor knowledge \cite{Piaget1952} is action-first:
infants learn to grasp, reach, and manipulate before they can represent
the physical laws that govern these actions. The propositions come later,
abstracted from accumulated action experience. Polanyi's tacit knowledge
\cite{Polanyi1966} --- ``we know more than we can tell'' --- includes
not only knowledge that has migrated \emph{below} propositional access
($K \to A$) but knowledge that was \emph{never propositional to begin
with} ($A$ without $K$).

\textbf{Bidirectionality.}
At no point in either direction does information cross a boundary from
``knowledge'' to ``action'' or vice versa. The encoding level determines
whether information appears propositional (high $c$, high $r$) or
procedural (low $c$, low $r$), but the information content is
continuous through the hierarchy. The entity tuple $E = (K, A, \sigma,
\Omega)$ is a notational convenience: $K$ is the high-encoding region of a
single structure, $A$ is the low-encoding region of the same structure,
and information flows in both directions between them as the entity
learns.
\end{proof}

\textbf{Implication.} The system does not ``decide'' whether information
belongs in $K$ or $A$. The encoding dynamics decide, and they operate
in both directions. Some information enters as explicit propositions and
compiles into action ($K \to A$: studying a recipe, then cooking without
reference to it). Some information enters as pure action and generates
propositions only on reflection ($A \to K$: pulling away from heat, then
learning why). Some information never becomes propositional at all: the
expert's ``feel'' for a material, the athlete's timing, the musician's
phrasing. These are $A$ without $K$ --- action knowledge that was never
derived from and may never produce explicit propositions.

The question ``where does this information go?'' has no answer because
there is no partition to sort into. There is only the encoding
hierarchy, and every piece of information has a position on it that
changes over time, entered from either direction.

This strengthens the inseparability result. The paper's theorem proves
the weak version: the \emph{content} of $K$ depends on $A$. The
continuum result proves the strong version: $K$ and $A$ are the same
structure viewed at different encoding depths, with bidirectional flow
between them. Separating them architecturally requires the system to
classify each piece of information as ``world model'' or ``action
policy'' at the moment of acquisition --- but for $A \to K$ information,
the classification is wrong at entry (it is pure action with no
propositional content), and for $K \to A$ information, the
classification becomes wrong over time (it was propositional but is now
compiled). The boundary is not merely blurry --- it is incoherent as a
classification criterion.
=======
By Corollary~\ref{cor:retention}(D), displacement occurs when a
candidate $q$ has strictly higher net marginal value than the current
occupant $p$. This is not information loss but information improvement:
the entity's net survival benefit improves with each displacement.
\end{proof}

\textbf{Remark (Free context).}
Define the \emph{free context} $C_{\mathrm{free}}(t) = C_n -
|\mathrm{ctx}_{\mathrm{unresolved}}(t)|$: the portion of bounded context
available for novel problems after existing knowledge maintenance.
Every proposition retained in $K$ reduces $C_{\mathrm{free}}$ by its
retrieval interference and edge maintenance overhead. Curation under
displacement (Corollary~\ref{cor:retention}(D)) is therefore a
$C_{\mathrm{free}}$-maximizing operation: displacing low-value
propositions frees context for higher-value candidates.

% ------------------------------------------------------------------
\section{Strong Inseparability: The Encoding Continuum}
\label{sec:continuum}
% ------------------------------------------------------------------

\begin{remark}[$K$ and $A$ Are Not Natural Kinds]
\label{rem:continuum}
Given Theorem~\ref{thm:normalization} ($K$ precipitates from $A$) and
the encoding hierarchy $E(c,r)$ of \cite{McCarthy2026b} (a continuous
space parameterized by runtime cost $c$ and reversibility $r$), $K$ and
$A$ are not disjoint categories but regions on $E(c,r)$. The encoding
hierarchy is empirically instantiated by the complementary learning
systems of McClelland et al.\ \cite{McClelland1995}: hippocampus (high
cost, high reversibility, fast episodic learning) and neocortex (low
cost, low reversibility, slow statistical consolidation). Memory
consolidation is movement along $E(c,r)$, operating in both directions.
\end{remark}

\textbf{Argument.}
We exhibit both directions of flow.

\textbf{$K \to A$ (declarative to procedural).}
A proposition $p$ grounding action $a$ begins at encoding level $L_n$
(active reasoning): explicitly represented, revisable. As the entity
repeatedly applies $a$ successfully, $p$ migrates to lower encoding
levels. At $L_{n-1}$ (learned heuristic), $p$ is retrievable but not
explicitly stated. At $L_1$ (developmental encoding), $p$ is compiled
into the execution of $a$. At $L_0$ (heritable encoding), $p$ is
implicit in the wiring that produces $a$. This is the
declarative-to-procedural transition observed in skill acquisition
\cite{FittsPosner1967}: the novice follows explicit rules ($K$ at
$L_n$); the expert executes embodied competence ($A$ at $L_1$).

\textbf{$A \to K$ (procedural to declarative).}
Novel information often enters as pure action. A child touches a hot
stove and pulls back. The first encoding is $A$: \emph{don't touch
that}. No proposition about temperature or tissue damage is represented.
Only later does the entity derive the propositional structure: ``the
stove is hot,'' ``hot things burn,'' ``heat transfers from higher to
lower temperature.'' Each proposition is derived from the action, not
the other way around.

\textbf{Bidirectionality.}
At no point in either direction does information cross a boundary from
``knowledge'' to ``action.'' The encoding level determines whether
information appears propositional (high $c$, high $r$) or procedural
(low $c$, low $r$), but the information content is continuous through the
hierarchy. The entity tuple $E = (K, A, \sigma, \Omega)$ is notational
convenience: $K$ is the high-encoding region of a single structure, $A$
is the low-encoding region, and information flows in both directions.

\textbf{Implication.} The system does not ``decide'' whether information
belongs in $K$ or $A$. The encoding dynamics decide. Some information
enters as explicit propositions and compiles into action ($K \to A$).
Some enters as pure action and generates propositions only on reflection
($A \to K$). Some never becomes propositional at all: the expert's
``feel'' for a material, the athlete's timing. These are $A$ without
$K$---action knowledge that was never derived from and may never produce
explicit propositions.

This strengthens the normalization result. The theorem proves the weak
version: $K$ separates from $A$ as the normalized shared propositions.
The continuum result proves the strong version: $K$ and $A$ are the same
structure viewed at different encoding depths, with bidirectional flow.
Separating them architecturally requires classifying each piece of
information as ``world model'' or ``action policy'' at acquisition, but
for $A \to K$ information, the classification is wrong at entry; for
$K \to A$ information, the classification becomes wrong over time.

\textbf{Remark on modularity of mind.} The strongest version of the
objection that $K$ and $A$ are natural kinds comes from Fodor's
modularity thesis \cite{Fodor1983}. Fodor argues that mental modules
are informationally encapsulated: they process input without access to
information held by other modules. Carruthers \cite{Carruthers2006}
extends the argument to central cognition.

The encapsulation is real. The axis of encapsulation is not. Fodor's
modules are encapsulated along the encoding hierarchy $E(c,r)$: a
processor at $L_1$ cannot access representations at $L_n$. The visual
module occupies a region of $E(c,r)$, not a side of a $K/A$ partition.
The same module holds both knowledge (edge statistics, color constancy
priors) and action (saccade targeting, vergence control) at its encoding
level. Encapsulation runs \emph{between encoding levels}, not between
knowledge and action. Physical modularity---hippocampus, neocortex,
motor cortex, cerebellum---is $E(c,r)$ physically instantiated, not
evidence for $K/A$ as a natural partition.
>>>>>>> 319f447daa46f8c1e949fb7929535d1aa6396025

% ------------------------------------------------------------------
\section{Objections and Scope}
\label{sec:objections}
% ------------------------------------------------------------------

<<<<<<< HEAD
\subsection{The unfalsifiable pruning criterion}

The pruning condition (P) requires $P(\exists\,\tau^* : B(p, \tau^*) >
0) = 0$ --- certain zero benefit across all reachable future
trajectories. No finite agent can verify this condition: for almost any
proposition, one can construct a scenario in which it is useful. Does
this make the pruning criterion operationally vacuous?

The unfalsifiability is a \emph{feature}, not a defect. It explains
three empirical phenomena:

\textbf{(1) Knowledge hoarding in the growth regime.} Biological
systems accumulate enormous quantities of seemingly useless information
before capacity forces curation. Children absorb facts, skills, and
associations indiscriminately. This is precisely the behavior the
theorem prescribes: in the growth regime ($|K| < K_{\max}$), the
unfalsifiable pruning criterion means \emph{never voluntarily discard}.
Hoarding is not a failure of selectivity --- it is the correct policy
under uncertain future benefit.

\textbf{(2) Forgetting is always forced, never voluntary.} Biological
forgetting is driven by capacity pressure (interference, decay,
displacement by new learning) rather than by a deliberate decision to
prune \cite{AndersonSchooler1991}. The theorem explains why:
the pruning condition cannot be verified, so pruning never occurs
voluntarily. Forgetting happens only in the curation regime, via
displacement (D), when a candidate with strictly higher marginal value
arrives.

\textbf{(3) The asymmetry operates within curation.} The displacement
criterion (D) is not symmetric ``higher value wins.'' The burden of
proof falls on displacement: the challenger must be \emph{strictly}
better ($\mathbb{E}[G(a_q, K)] > \mathbb{E}[G(a_p, K)]$, not $\geq$).
Ties go to the incumbent. This is the retention asymmetry operating
within the curation regime --- not merely ``keep things when you have
room'' but ``default to retention even when choosing what to displace.''

\subsection{The scope of survival pressure}

The retention asymmetry depends on $C_{\text{absent}}$ dominating all
finite costs (Lemma~\ref{lem:absence}), which requires survival
pressure. But most knowledge does not appear survival-critical. The
capital of France does not ground a survival-relevant action. Does the
theorem apply only to a small subset of $K$?

The answer turns on what the entity can know \emph{in advance} about
which propositions are survival-critical. Under genuine uncertainty
about future domains of encounter, the entity \emph{cannot} partition
its knowledge into ``survival-relevant'' and ``not.'' If it could, it
would already have the information the retention criterion is trying to
preserve: which propositions ground which actions in which domains.

Under uncertainty about the distribution of future encounters:

\begin{itemize}
\item Any proposition \emph{might} ground a survival-critical action in
   an unforeseen domain. The proposition ``Paris is the capital of
   France'' is survival-irrelevant on an ordinary Tuesday and
   survival-critical when navigating an unfamiliar city to reach a
   hospital.
\item The cost of wrongly classifying a proposition as non-critical
   ($C_{\text{absent}}$) dominates the cost of wrongly classifying it as
   critical ($c_{\text{retain}}$), by the same asymmetry the theorem
   establishes.
\item Therefore the conservative criterion extends to \emph{all}
   propositions whose future relevance is uncertain --- which, for any
   finite agent with incomplete knowledge of its future trajectory, is
   nearly all propositions.
\end{itemize}

This is why humans retain vast stores of ``useless'' knowledge. The
retention is not irrational hoarding but a rational response to the
impossibility of verifying irrelevance under uncertain future
encounters. The scope of survival pressure is wider than it appears
because the entity cannot narrow it without the very knowledge the
retention criterion governs.

\textbf{Remark.} For entities genuinely not under survival pressure ---
archival databases, reference libraries, static knowledge stores --- the
retention asymmetry weakens. $C_{\text{absent}}$ is finite (the cost of
a missing record is a failed query, not non-survival), and the symmetric
marginal value comparison of the curation regime becomes the primary
decision rule. The theorem's scope is explicitly survival-pressured
entities, which includes all biological organisms and any artificial
agent whose continued operation depends on performance.
=======
\subsection{The overlap assumption}

The normalization theorem requires $\beta > 0$: at least some
propositions are shared across action policies. Is this assumption
restrictive?

In practice, $\beta > 0$ is the norm for any entity with a
non-trivial action repertoire. Physical laws, spatial relations,
object identities, social norms, and causal regularities ground
multiple action policies simultaneously. An entity whose action
policies share \emph{no} propositions would have a purely reflexive
repertoire---each action depending only on $W_{\text{obs}}$, with
no propositional grounding at all. Such an entity has $K = \emptyset$
and the normalization question does not arise. For any entity with
$K \neq \emptyset$ and $|A| > 1$, overlap is the expected condition.

\subsection{The scope of survival pressure}

The retention asymmetry (Corollary~\ref{cor:retention}) depends on
$C_{\text{absent}}$ dominating all finite costs, which requires
survival pressure. But most knowledge does not appear
survival-critical.

The answer turns on what the entity can know \emph{in advance} about
which propositions are survival-critical. Under genuine uncertainty
about future domains of encounter, the entity cannot partition its
knowledge into ``survival-relevant'' and ``not.'' Any proposition
might ground a survival-critical action in an unforeseen domain. The
proposition ``Paris is the capital of France'' is survival-irrelevant
on an ordinary Tuesday and survival-critical when navigating an
unfamiliar city to reach a hospital.

The cost of wrongly classifying a proposition as non-critical
($C_{\text{absent}}$) dominates the cost of wrongly classifying it
as critical ($c_{\text{retain}}$), by the same asymmetry the
corollary establishes. The conservative criterion therefore extends
to all propositions whose future relevance is uncertain---which, for
any finite entity, is nearly all propositions.

For entities genuinely not under survival pressure---archival databases,
static knowledge stores---the retention asymmetry weakens.
$C_{\text{absent}}$ is finite, and the symmetric marginal value
comparison of the curation regime becomes the primary decision rule.

\subsection{The unfalsifiable pruning criterion}

The pruning condition (P) requires $P(\exists\,\tau^*) = 0$: certain
zero benefit across all reachable trajectories. No finite agent can
verify this. Does this make pruning operationally vacuous?

The unfalsifiability is a feature, not a defect. It explains three
empirical phenomena:

\textbf{(1) Knowledge hoarding in the growth regime.} Biological
systems accumulate enormous quantities of seemingly useless information
before capacity forces curation. This is the behavior
Corollary~\ref{cor:retention} prescribes: in the growth regime, the
unfalsifiable pruning criterion means \emph{never voluntarily discard}.

\textbf{(2) Forgetting is always forced, never voluntary.} Biological
forgetting is driven by capacity pressure rather than deliberate
pruning \cite{AndersonSchooler1991}. The corollary explains why:
the pruning condition cannot be verified, so forgetting happens only
via displacement (D).

\textbf{(3) The asymmetry operates within curation.} The displacement
criterion is not symmetric ``higher value wins.'' The burden of proof
falls on displacement: the challenger must be strictly better. Ties go
to the incumbent.
>>>>>>> 319f447daa46f8c1e949fb7929535d1aa6396025

% ------------------------------------------------------------------
\section{Engagement with Prior Work}
\label{sec:discussion}
% ------------------------------------------------------------------

<<<<<<< HEAD
\textbf{Dupoux, LeCun, and Malik \cite{DupouxLeCunMalik2026}.}
Their proposed architecture separates System~A (observation and world
modeling) from System~B (action and planning), with each system having
its own training objective and a meta-controller routing tasks between
them. A natural objection to our result is that implementation modularity
--- separate modules with an interface --- is compatible with
content-dependence. A database indexed by queries has content shaped by
the query workload, but the database and query engine are separate
modules. Does $K/A$ inseparability prohibit modular implementation?

No. The claim is about \emph{optimization}, not \emph{implementation}.
Separate modules are compatible with the theorem provided they share a
training objective that references both $K$ and $A$: A's loss must
include action relevance, B's loss must include knowledge availability.
``Separate modules with a shared loss'' is co-optimization with a
modular implementation --- already inseparability at the optimization
level.

What Theorem~\ref{thm:inseparability} prohibits is what Dupoux, LeCun,
and Malik specifically propose: \emph{separate training objectives} for
A and B. Their System~A is trained on prediction (world modeling);
their System~B is trained on planning (action optimization). These are
independent objectives. Under survival pressure, independent
optimization of A produces a world model that retains propositions
irrelevant to B's action space (wasting bounded capacity) or fails to
retain propositions critical to B (risking
$U > U_{\text{lethal}}^d$). The content of A must reference B's
capabilities, and the content of B must reference A's retained
propositions. Separate modules, shared loss: this is the structural
constraint.

\textbf{Information Bottleneck \cite{Tishby2000, Tishby2011}.}
The information bottleneck formalizes the trade-off between compression
and relevance: retain information about input $X$ that is relevant to
output $Y$, subject to a rate constraint. The retention threshold is
controlled by the Lagrange multiplier $\beta$, which trades off
compression against prediction. In practice, $\beta$ is a
hyperparameter chosen by the system designer. Our result recovers a
version of this trade-off but \emph{derives} the threshold rather than
positing it. The threshold is $C_{\text{absent}}$: the cost of a
proposition's absence when needed. Under survival pressure,
$C_{\text{absent}}$ is bounded below by entity non-survival
(Lemma~\ref{lem:absence}), which dominates all finite costs. The
resulting retention criterion is therefore maximally conservative ---
retain whenever benefit is uncertain --- not because of a design choice
for large $\beta$ but because the problem structure demands it.

\textbf{Gibson \cite{Gibson1979}.}
Gibson's affordances are properties of the environment specified relative
to an organism's action capabilities: a surface ``affords'' walking for
a creature with legs but not for a creature without. This is knowledge
indexed by action under a different name. Corollary~\ref{cor:indexed}
formalizes Gibson's insight: two entities with the same perceptual
apparatus ($\sigma$) but different action capabilities ($A$) maintain
different knowledge ($K$). What Theorem~\ref{thm:inseparability} adds
to Gibson is a formal selection criterion. Gibson describes affordances
as directly perceived but does not specify which affordances are
retained when capacity is bounded. The asymmetric retention criterion
--- retain when benefit is uncertain, prune only when certainly zero ---
provides the missing selection rule.

\textbf{Anderson and Schooler \cite{AndersonSchooler1991}.}
Anderson and Schooler demonstrated that the probability of needing a
memory item mirrors the statistical structure of the environment:
items encountered more recently and more frequently are more likely to
be needed. This is an empirical prediction of
Theorem~\ref{thm:inseparability} applied to biological $K$. If
retention is indexed by action relevance, and action relevance tracks
environmental statistics (more frequent situations require more
frequent action), then memory retention curves should mirror
environmental frequency distributions --- which is precisely what
Anderson and Schooler observe.

\textbf{Simon \cite{Simon1955}.}
Simon's bounded rationality establishes that agents with limited
cognitive resources satisfice rather than optimize: they seek solutions
that are good enough rather than searching for the best. The retention
criterion of Theorem~\ref{thm:inseparability} is a satisficing
criterion in Simon's sense. The entity does not compute the optimal
knowledge set (which would require evaluating all possible subsets of
propositions --- a combinatorial problem). Instead, it applies a simple
rule: retain if benefit is not certainly zero, displace if a better
candidate arrives. This rule is computationally tractable and, by the
submodularity guarantee of the greedy criterion, achieves a constant-factor
approximation to optimal.
=======
\textbf{Codd \cite{Codd1970}.}
The normalization theorem is the knowledge-system analogue of Codd's
relational normalization. Codd identified that redundant storage of
data across relations produces update anomalies and wasted capacity;
normal forms extract shared dependencies into referenced tables. We
prove that the same pressure operates on bounded cognitive systems:
shared propositions across action policies produce redundant context
load and update inconsistency; $K$ is the normal form that eliminates
both. The contribution is showing that normalization is not merely a
good engineering practice but a necessary consequence of bounded context.

\textbf{Dupoux, LeCun, and Malik \cite{DupouxLeCunMalik2026}.}
Their proposed architecture separates System~A (observation and world
modeling) from System~B (action and planning), with separate training
objectives. Corollary~\ref{cor:inseparability} shows that separate
training objectives violate a structural constraint: the content of
the world model must reference the action space (which propositions are
action-relevant), and the action space must reference the world model
(which actions are groundable). Separate modules with a shared loss is
compatible; separate training objectives is not.

\textbf{Information Bottleneck \cite{Tishby2000, Tishby2011}.}
The information bottleneck retains information relevant to an output
variable subject to a rate constraint, with threshold controlled by
Lagrange multiplier $\beta$. Our retention criterion
(Corollary~\ref{cor:retention}) recovers a version of this trade-off
but derives the threshold rather than positing it: under survival
pressure, $C_{\text{absent}}$ is bounded below by entity non-survival,
producing a maximally conservative retention rule. The normalization
theorem adds a structural result the information bottleneck lacks:
not just \emph{which} information to retain, but \emph{how} to
organize it (shared propositions separated from policies).

\textbf{Active inference \cite{Friston2010, ParrPezzuloFriston2022}.}
Active inference argues that perception and action are jointly optimized
under variational free energy minimization. This is the closest existing
framework to $K/A$ inseparability. What
Theorem~\ref{thm:normalization} adds is the organizational constraint
under hard capacity limits. When context is unbounded, all information
can be retained and the coupling is a soft preference. Under bounded
$C_n$, the coupling becomes a hard constraint: propositions must be
displaced, and the displacement criterion must reference the action
space. Active inference tells the entity what to do; the normalization
theorem tells it how to organize what it knows while doing it.

\textbf{Gibson \cite{Gibson1979}.}
Gibson's affordances---properties of the environment specified relative to
an organism's action capabilities---are knowledge indexed by action under
a different name. Corollary~\ref{cor:inseparability} formalizes Gibson's
insight. What the normalization theorem adds is a structural account:
affordances are not merely perceived but \emph{factored out} from action
policies as shared propositions, and this factoring is forced by bounded
capacity.

\textbf{Anderson and Schooler \cite{AndersonSchooler1991}.}
They demonstrated that memory retention statistics mirror environmental
statistics. This is an empirical prediction of
Corollary~\ref{cor:retention}: if retention is indexed by action
relevance, and action relevance tracks environmental statistics, then
retention curves should mirror environmental frequency
distributions---precisely what they observe.

\textbf{Simon \cite{Simon1955}.}
The retention criterion is a satisficing criterion in Simon's sense.
The entity does not compute the optimal knowledge set (a combinatorial
problem) but applies a simple rule: retain if benefit is not certainly
zero, displace if a better candidate arrives. When the benefit function
is submodular, the greedy criterion achieves a $(1 - 1/e)$-approximation
to optimal \cite{Nemhauser1978}. In the general case, the guarantee does
not apply, but the criterion remains a well-defined satisficing
heuristic.
>>>>>>> 319f447daa46f8c1e949fb7929535d1aa6396025

% ------------------------------------------------------------------
\section{Conclusion}
\label{sec:conclusion}
% ------------------------------------------------------------------

<<<<<<< HEAD
The inseparability of knowledge and action under survival pressure
operates at two levels.

The \emph{weak} inseparability (Theorem~\ref{thm:inseparability}): the
content of $K$ is constitutively determined by $A$. The retention
criterion references the action space at every step; the asymmetry
between the cost of retention (finite) and the cost of absence
(potentially unbounded under survival pressure) produces a maximally
conservative retention rule indexed by action relevance. Systems may
implement $K$ and $A$ as separate modules, but cannot give them
independent training objectives without accepting information loss or
reconstruction cost.

The \emph{strong} inseparability
(Corollary~\ref{cor:continuum}): $K$ and $A$ are not natural kinds but
regions on the encoding hierarchy $E(c,r)$, with information flowing in
both directions. Some information enters as explicit knowledge and
compiles into action ($K \to A$: studying a recipe, then cooking without
reference). Some enters as pure action and generates propositions only
on reflection ($A \to K$: pulling away from a hot stove, then learning
that hot things burn). Some never becomes propositional at all. The
system does not decide where information belongs because there is no
partition to sort into --- only a continuum parameterized by runtime
cost and reversibility, entered from either direction.

The strong result sharpens the critique of observation-action
separation. The problem is not merely that separate training objectives
produce suboptimal retention (the weak claim). The problem is that the
separation requires classifying each piece of information as ``world
model'' or ``action policy'' at the moment of acquisition --- but for
$A \to K$ information (the dominant mode of early learning), the
classification is wrong at entry, and for $K \to A$ information, the
classification becomes wrong over time. The boundary is not merely
blurry --- it is incoherent as a classification criterion.

The knowledge graph $K$ precipitates from the action space $A$ as the
normalization of shared propositions --- converting multiplicative
scaling ($O(n \cdot k)$) to additive scaling ($O(n + k)$) --- and this
precipitation is forced by bounded context, not chosen by a designer.
The graph is not a store of declarative facts separate from procedural
skills; it is the shared structure underlying both, factored out to
avoid redundancy under bounded capacity.
=======
Knowledge as normalization operates at two levels.

The \emph{structural} level (Theorem~\ref{thm:normalization}):
bounded active context with overlapping action policies forces the
extraction of shared propositions into a separated structure $K$,
converting multiplicative context load ($O(n \cdot k)$) to additive
load ($O(n + k')$) and eliminating update anomalies. This is not a
design choice but a scaling constraint. The entity that does not
normalize cannot scale its action repertoire beyond $C_n / k$ policies.

The \emph{content} level (Corollaries~\ref{cor:retention}--\ref{cor:inseparability}):
the propositions retained in $K$ are constitutively determined by
$A$. The retention criterion is asymmetric (retain when benefit is
uncertain; prune only when certainly zero) and derived from the cost
structure of bounded context under survival pressure, not from a
design parameter. Systems may implement $K$ and $A$ as separate
modules, but cannot give them independent training objectives without
accepting information loss or reconstruction cost.

The \emph{strong} inseparability (Remark~\ref{rem:continuum}): $K$
and $A$ are not natural kinds but regions on the encoding hierarchy
$E(c,r)$, with information flowing in both directions. The boundary
is not merely blurry---it is incoherent as a classification criterion.

The knowledge graph precipitates from the action space as the
normalization of shared propositions. It is not a store of declarative
facts separate from procedural skills; it is the shared structure
underlying both, factored out to avoid redundancy under bounded
capacity.

\section*{Acknowledgments}

The author used Claude (Anthropic) for drafting assistance, literature review, and \LaTeX{} preparation. All theoretical content, proofs, and arguments are the author's own.
>>>>>>> 319f447daa46f8c1e949fb7929535d1aa6396025

\bibliographystyle{plain}
\begin{thebibliography}{99}

\bibitem{Codd1970}
E.~F. Codd.
\newblock A relational model of data for large shared data banks.
\newblock \emph{Communications of the ACM}, 13(6):377--387, 1970.

\bibitem{AndersonSchooler1991}
J.~R. Anderson and L.~J. Schooler.
\newblock Reflections of the environment in memory.
\newblock \emph{Psychological Science}, 2(6):396--408, 1991.

\bibitem{Cover2006}
T.~M. Cover and J.~A. Thomas.
\newblock \emph{Elements of Information Theory}.
\newblock Wiley-Interscience, 2nd edition, 2006.

<<<<<<< HEAD
\bibitem{Dreyfus1986}
H.~L. Dreyfus and S.~E. Dreyfus.
\newblock \emph{Mind Over Machine: The Power of Human Intuition and Expertise
  in the Era of the Computer}.
\newblock Free Press, 1986.
=======
\bibitem{Carruthers2006}
P.~Carruthers.
\newblock \emph{The Architecture of the Mind: Massive Modularity and the
  Flexibility of Thought}.
\newblock Oxford University Press, 2006.
>>>>>>> 319f447daa46f8c1e949fb7929535d1aa6396025

\bibitem{DixitPindyck1994}
A.~K. Dixit and R.~S. Pindyck.
\newblock \emph{Investment Under Uncertainty}.
\newblock Princeton University Press, 1994.

\bibitem{DupouxLeCunMalik2026}
E.~Dupoux, Y.~LeCun, and J.~Malik.
\newblock Why {AI} systems don't learn and what to do about it: Lessons on
  autonomous learning from cognitive science.
\newblock \emph{arXiv preprint arXiv:2603.15381}, 2026.

<<<<<<< HEAD
\bibitem{FittsPostner1967}
=======
\bibitem{Fodor1983}
J.~A. Fodor.
\newblock \emph{The Modularity of Mind}.
\newblock MIT Press, 1983.

\bibitem{Friston2010}
K.~Friston.
\newblock The free-energy principle: A unified brain theory?
\newblock \emph{Nature Reviews Neuroscience}, 11(2):127--138, 2010.

\bibitem{ParrPezzuloFriston2022}
T.~Parr, G.~Pezzulo, and K.~J. Friston.
\newblock \emph{Active Inference: The Free Energy Principle in Mind, Brain, and
  Behavior}.
\newblock MIT Press, 2022.

\bibitem{FittsPosner1967}
>>>>>>> 319f447daa46f8c1e949fb7929535d1aa6396025
P.~M. Fitts and M.~I. Posner.
\newblock \emph{Human Performance}.
\newblock Brooks/Cole, 1967.

\bibitem{Gibson1979}
J.~J. Gibson.
\newblock \emph{The Ecological Approach to Visual Perception}.
\newblock Houghton Mifflin, 1979.

\bibitem{McClelland1995}
J.~L. McClelland, B.~L. McNaughton, and R.~C. O'Reilly.
\newblock Why there are complementary learning systems in the hippocampus and
  neocortex: Insights from the successes and failures of connectionist models
  of learning and memory.
\newblock \emph{Psychological Review}, 102(3):419--457, 1995.

\bibitem{McCarthy2026a}
P.~D. McCarthy.
\newblock Graph structure is necessary for information preservation under
  bounded context.
\newblock Technical report, 2026.

\bibitem{McCarthy2026b}
P.~D. McCarthy.
<<<<<<< HEAD
\newblock Escalation dominates top-down allocation under bounded context.
\newblock Working paper, 2026.

\bibitem{Miller1956}
G.~A. Miller.
\newblock The magical number seven, plus or minus two: Some limits on our
  capacity for processing information.
\newblock \emph{Psychological Review}, 63(2):81--97, 1956.
=======
\newblock Convergence of control architectures to escalation under bounded context.
\newblock Working paper, 2026.

\bibitem{McCarthy2026d}
P.~D. McCarthy.
\newblock Evaluation-driven descent, encoding permanence, and the structural
  invariants of intelligence.
\newblock Manuscript, 2026.
>>>>>>> 319f447daa46f8c1e949fb7929535d1aa6396025

\bibitem{Nemhauser1978}
G.~L. Nemhauser, L.~A. Wolsey, and M.~L. Fisher.
\newblock An analysis of approximations for maximizing submodular set
  functions --- {I}.
\newblock \emph{Mathematical Programming}, 14(1):265--294, 1978.

<<<<<<< HEAD
\bibitem{Polanyi1966}
M.~Polanyi.
\newblock \emph{The Tacit Dimension}.
\newblock Doubleday, 1966.

\bibitem{Piaget1952}
J.~Piaget.
\newblock \emph{The Origins of Intelligence in Children}.
\newblock International Universities Press, 1952.

\bibitem{Pearl1988}
J.~Pearl.
\newblock \emph{Probabilistic Reasoning in Intelligent Systems: Networks of
  Plausible Inference}.
\newblock Morgan Kaufmann, 1988.

=======
>>>>>>> 319f447daa46f8c1e949fb7929535d1aa6396025
\bibitem{Shannon1948}
C.~E. Shannon.
\newblock A mathematical theory of communication.
\newblock \emph{Bell System Technical Journal}, 27(3):379--423, 1948.

\bibitem{Simon1955}
H.~A. Simon.
\newblock A behavioral model of rational choice.
\newblock \emph{The Quarterly Journal of Economics}, 69(1):99--118, 1955.

\bibitem{Tishby2000}
N.~Tishby, F.~C. Pereira, and W.~Bialek.
\newblock The information bottleneck method.
\newblock In \emph{Proceedings of the 37th Annual Allerton Conference on
  Communication, Control, and Computing}, pages 368--377, 2000.

\bibitem{Tishby2011}
N.~Tishby and N.~Zaslavsky.
\newblock Deep learning and the information bottleneck principle.
\newblock In \emph{IEEE Information Theory Workshop}, pages 1--5, 2015.

\end{thebibliography}

\end{document}
